\documentclass[aspectratio=169,10pt]{beamer}

\usepackage[spanish,es-noquoting]{babel}
\usepackage[utf8]{inputenc}
\usepackage[T1]{fontenc}
\usepackage{lmodern}
\usepackage{microtype}
\usepackage{tikz}
\usepackage{pgfplots}
\usepackage{booktabs}
\usepackage{amsmath}
\usetikzlibrary{arrows.meta,positioning,calc,shapes.geometric,fit,decorations.pathreplacing}
\pgfplotsset{compat=1.18}

\definecolor{DeepBlue}{HTML}{1F4E79}
\definecolor{Accent}{HTML}{E36C0A}
\definecolor{SoftBlue}{HTML}{EAF2F8}
\definecolor{SoftOrange}{HTML}{FFF1E5}
\definecolor{GrayText}{HTML}{3A3A3A}
\definecolor{DarkGray}{HTML}{555555}
\definecolor{Good}{HTML}{2E7D32}
\definecolor{Warn}{HTML}{B35C00}

\setbeamercolor{normal text}{fg=GrayText,bg=white}
\setbeamercolor{frametitle}{fg=DeepBlue,bg=white}
\setbeamercolor{title}{fg=DeepBlue}
\setbeamercolor{subtitle}{fg=DarkGray}
\setbeamercolor{structure}{fg=DeepBlue}
\setbeamertemplate{navigation symbols}{}
\setbeamertemplate{frametitle}{\vspace{0.25cm}\insertframetitle\par\vspace{0.1cm}\hrule height 0.5pt\vspace{0.2cm}}
\setbeamertemplate{footline}{\hfill\scriptsize\insertframenumber/10\hspace{0.35cm}\vspace{0.15cm}}

\newcommand{\tagbox}[2]{\tikz[baseline=(X.base)]\node[rounded corners=2pt,fill=#1,inner xsep=5pt,inner ysep=2pt](X){\scriptsize\textbf{#2}};}
\newcommand{\source}[1]{\vfill{\tiny\color{DarkGray}Fuente: #1}}

\title{Crítica a \emph{Performance Evaluation for the Hash Generation Phase of a Democratic Blockchain}}
\subtitle{Eficiencia computacional, consenso y sostenibilidad}
\author{Andrés Arias, Catalina Jaramillo y Rafael Marulanda}
\institute{Pontificia Universidad Javeriana}
\date{4 de septiembre de 2026}

\begin{document}

% 1
\begin{frame}[plain]
  \vspace{0.4cm}
  {\Large\bfseries\color{DeepBlue} Crítica a Lugo y Pedraza (2020)}\\[0.15cm]
  {\large \emph{Performance Evaluation for the Hash Generation Phase of a Democratic Blockchain}}\\[0.25cm]
  {\small Andrés Arias, Catalina Jaramillo y Rafael Marulanda}\\[0.06cm]
  {\scriptsize Pontificia Universidad Javeriana}\\[0.35cm]
  \begin{tikzpicture}[x=1cm,y=1cm]
    \node[draw=DeepBlue, rounded corners=5pt, line width=1pt, fill=SoftBlue, minimum width=3.4cm, minimum height=1.0cm] (a) at (0,0) {\textbf{Hashing}};
    \node[draw=Accent, rounded corners=5pt, line width=1pt, fill=SoftOrange, minimum width=3.4cm, minimum height=1.0cm] (b) at (4.2,0) {\textbf{Hardware}};
    \node[draw=DarkGray, rounded corners=5pt, line width=1pt, fill=white, minimum width=3.4cm, minimum height=1.0cm] (c) at (8.4,0) {\textbf{Consenso}};
    \draw[-{Latex[length=2.5mm]},thick,DeepBlue] (a) -- (b);
    \draw[-{Latex[length=2.5mm]},thick,Accent] (b) -- (c);
    \node[font=\small,align=center,text width=10.5cm] at (4.2,-1.1) {El artículo mide bien una subrutina computacional, pero no prueba superioridad sistémica del protocolo democrático.};
  \end{tikzpicture}
  \vfill
  {\footnotesize Nakamoto (2008); Paul, Sarkar y Mukherjee (2014); Lugo y Pedraza (2020).}
\end{frame}

% 2
\begin{frame}{Pregunta de la revisión}
\begin{columns}[T,totalwidth=\textwidth]
\column{0.53\textwidth}
\textbf{Objeto del artículo}\vspace{0.15cm}
\begin{itemize}\setlength\itemsep{0.25cm}
  \item Fase de generación de hashes de una blockchain democrática.
  \item Comparación de CPU, GPU y clúster.
  \item Medición de tiempo y estimación de energía.
\end{itemize}

\column{0.44\textwidth}
\begin{tikzpicture}
  \node[rounded corners=5pt, fill=SoftBlue, draw=DeepBlue, text width=4.8cm, inner sep=8pt] {
  \textbf{Tesis crítica}\par\vspace{0.1cm}
  La eficiencia por hash es relevante, pero no basta para inferir seguridad, descentralización, escalabilidad o sostenibilidad del sistema completo.
  };
\end{tikzpicture}
\end{columns}
\end{frame}

% 3
\begin{frame}{Bitcoin PoW vs. minería democrática}
\centering
\begin{tikzpicture}[node distance=0.75cm, every node/.style={font=\small}]
  \node[draw=DeepBlue, fill=SoftBlue, rounded corners, minimum width=4.7cm, minimum height=0.8cm] (pow1) {Bitcoin PoW};
  \node[below=of pow1, draw=DeepBlue, rounded corners, text width=4.7cm, align=center, minimum height=1.0cm] (pow2) {Encontrar primero un hash \(H<T\)};
  \node[below=of pow2, draw=DeepBlue, rounded corners, text width=4.7cm, align=center, minimum height=1.0cm] (pow3) {Probabilidad proporcional al poder de cómputo};
  \node[below=of pow3, draw=DeepBlue, rounded corners, text width=4.7cm, align=center, minimum height=1.0cm] (pow4) {Seguridad por trabajo acumulado};

  \node[right=1.3cm of pow1, draw=Accent, fill=SoftOrange, rounded corners, minimum width=4.7cm, minimum height=0.8cm] (dem1) {Protocolo democrático};
  \node[below=of dem1, draw=Accent, rounded corners, text width=4.7cm, align=center, minimum height=1.0cm] (dem2) {Generar hashes en una ventana temporal};
  \node[below=of dem2, draw=Accent, rounded corners, text width=4.7cm, align=center, minimum height=1.0cm] (dem3) {Seleccionar el menor hash};
  \node[below=of dem3, draw=Accent, rounded corners, text width=4.7cm, align=center, minimum height=1.0cm] (dem4) {Menos carrera local, más coordinación};

  \draw[-{Latex[length=2mm]}, thick] (pow1) -- (pow2);
  \draw[-{Latex[length=2mm]}, thick] (pow2) -- (pow3);
  \draw[-{Latex[length=2mm]}, thick] (pow3) -- (pow4);
  \draw[-{Latex[length=2mm]}, thick] (dem1) -- (dem2);
  \draw[-{Latex[length=2mm]}, thick] (dem2) -- (dem3);
  \draw[-{Latex[length=2mm]}, thick] (dem3) -- (dem4);
\end{tikzpicture}
\end{frame}

% 4
\begin{frame}{Alcance real del experimento}
\centering
\begin{tikzpicture}[node distance=1.0cm, every node/.style={font=\small}]
  \node[draw=Good, fill=Good!10, rounded corners, minimum width=3.5cm, minimum height=1cm] (gen) {\textbf{Generación}};
  \node[right=1.1cm of gen, draw=Warn, fill=Warn!10, rounded corners, minimum width=3.5cm, minimum height=1cm] (dif) {Difusión};
  \node[right=1.1cm of dif, draw=Warn, fill=Warn!10, rounded corners, minimum width=3.5cm, minimum height=1cm] (val) {Validación};
  \draw[-{Latex[length=2.5mm]}, thick] (gen) -- (dif);
  \draw[-{Latex[length=2.5mm]}, thick] (dif) -- (val);
  \node[below=0.55cm of gen, text width=3.5cm, align=center] {Fase evaluada};
  \node[below=0.55cm of dif, text width=3.5cm, align=center] {No medida};
  \node[below=0.55cm of val, text width=3.5cm, align=center] {No medida};
  \draw[decorate,decoration={brace,amplitude=5pt,mirror},thick,Warn] ($(dif.south west)+(0,-0.95)$) -- ($(val.south east)+(0,-0.95)$) node[midway,below=0.18cm,text width=6.5cm,align=center] {Costos sistémicos: mensajes, latencia, escalabilidad, consenso};
\end{tikzpicture}
\vspace{0.25cm}
\begin{center}
\tagbox{Good!15}{Punto a favor}: analiza la fase computacionalmente más demandante.\quad
\tagbox{Warn!15}{Límite}: no evalúa la red completa.
\end{center}
\end{frame}

% 5
\begin{frame}{Diseño experimental}
\begin{columns}[T,totalwidth=\textwidth]
\column{0.48\textwidth}
\begin{tikzpicture}[every node/.style={font=\small}, node distance=0.55cm]
  \node[draw=DeepBlue, rounded corners, fill=SoftBlue, minimum width=4.8cm, minimum height=0.8cm] (h) {Carga común};
  \node[below=of h, draw=DeepBlue, rounded corners, text width=4.8cm, align=center, inner sep=6pt] (b) {Bloque Bitcoin \#453194\\nonce hasta \(4\,003\,888\,472\)\\doble SHA-256};
  \node[below=of b, draw=DeepBlue, rounded corners, text width=4.8cm, align=center, inner sep=6pt] (d) {Descomposición cíclica\\iteraciones repartidas por hilo};
  \draw[-{Latex[length=2mm]},thick] (h) -- (b);
  \draw[-{Latex[length=2mm]},thick] (b) -- (d);
\end{tikzpicture}

\column{0.49\textwidth}
\begin{tikzpicture}[every node/.style={font=\small}, node distance=0.4cm]
  \node[draw=DarkGray, rounded corners, text width=5cm, inner sep=5pt] (cpu) {\textbf{CPU}: Intel i7-7700HQ\\POSIX, OpenMP};
  \node[below=of cpu, draw=DarkGray, rounded corners, text width=5cm, inner sep=5pt] (gpu) {\textbf{GPU}: GTX 1050\\CUDA, OpenCL};
  \node[below=of gpu, draw=DarkGray, rounded corners, text width=5cm, inner sep=5pt] (mpi) {\textbf{Clúster}: 4 instancias GCP\\Open MPI};
  \node[below=0.7cm of mpi, fill=SoftOrange, draw=Accent, rounded corners, text width=5cm, inner sep=6pt] {No incluye ASIC: limita la comparación con minería competitiva de Bitcoin.};
\end{tikzpicture}
\end{columns}
\end{frame}

% 6
\begin{frame}{Resultado principal: tiempo de ejecución}
\centering
\begin{tikzpicture}
\begin{axis}[
  width=11.3cm,height=5.5cm,
  ybar,
  bar width=22pt,
  ymin=0,
  ylabel={Segundos},
  symbolic x coords={CPU,CUDA,OpenCL},
  xtick=data,
  nodes near coords,
  nodes near coords align={vertical},
  every node near coord/.append style={font=\small},
  axis lines*=left,
  enlarge x limits=0.3,
  ymajorgrids=true,
  grid style={gray!20},
  tick label style={font=\small},
  label style={font=\small},
  fill=DeepBlue!65,
]
\addplot coordinates {(CPU,1351) (CUDA,120) (OpenCL,106)};
\end{axis}
\end{tikzpicture}

\vspace{0.1cm}
\small Las GPU reducen el tiempo aproximadamente un orden de magnitud frente a la mejor configuración CPU.
\end{frame}

% 7
\begin{frame}{Resultado energético estimado}
\centering
\begin{tikzpicture}
\begin{axis}[
  width=11.5cm,height=5.5cm,
  ybar,
  bar width=18pt,
  ymin=0,
  ymax=110,
  ylabel={Wh estimados},
  symbolic x coords={POSIX,OpenMP,CUDA,Open MPI},
  xtick=data,
  nodes near coords,
  nodes near coords align={vertical},
  every node near coord/.append style={font=\scriptsize},
  axis lines*=left,
  enlarge x limits=0.22,
  ymajorgrids=true,
  grid style={gray!20},
  tick label style={font=\small},
  label style={font=\small},
  fill=Accent!70,
]
\addplot coordinates {(POSIX,21.05) (OpenMP,21.10) (CUDA,4.66) (Open MPI,100.20)};
\end{axis}
\end{tikzpicture}

\small La comparación distingue potencia instantánea de energía total: terminar antes puede consumir menos.
\end{frame}

% 8
\begin{frame}{Argumentos a favor}
\begin{tikzpicture}[every node/.style={font=\small}]
  \node[draw=Good, fill=Good!10, rounded corners, text width=5.8cm, inner sep=7pt] (a) at (0,2.0) {\textbf{Delimitación clara}\\Encabezado, nonce, función hash y plataformas quedan definidos.};
  \node[draw=Good, fill=Good!10, rounded corners, text width=5.8cm, inner sep=7pt] (b) at (6.6,2.0) {\textbf{Comparación útil}\\CPU, GPU y clúster bajo la misma carga de trabajo.};
  \node[draw=Good, fill=Good!10, rounded corners, text width=5.8cm, inner sep=7pt] (c) at (0,0.35) {\textbf{Resultado robusto}\\CUDA y OpenCL coinciden en la ventaja de GPU para hashing masivo.};
  \node[draw=Good, fill=Good!10, rounded corners, text width=5.8cm, inner sep=7pt] (d) at (6.6,0.35) {\textbf{Dimensión energética}\\Incluye energía estimada, no solo velocidad de ejecución.};
  \node[draw=DeepBlue, fill=SoftBlue, rounded corners, text width=12.4cm, inner sep=8pt] at (3.3,-1.45) {\textbf{Aporte central:} la arquitectura de hardware condiciona el costo computacional y energético de ejecutar la fase de generación de hashes.};
\end{tikzpicture}
\end{frame}

% 9
\begin{frame}{Argumentos en contra}
\centering
\begin{tikzpicture}[every node/.style={font=\small}, node distance=0.8cm]
  \node[draw=Warn, fill=Warn!10, rounded corners, text width=4.0cm, align=center, inner sep=7pt] (local) {Mejora local\\\textbf{hash generation}};
  \node[right=1.1cm of local, draw=DarkGray, rounded corners, text width=4.0cm, align=center, inner sep=7pt] (system) {Conclusión sistémica\\seguridad, escala, sostenibilidad};
  \draw[-{Latex[length=2.5mm]}, very thick, dashed, Warn] (local) -- node[above]{\scriptsize salto no demostrado} (system);

  \node[below=1.0cm of local, draw=Warn, fill=SoftOrange, rounded corners, text width=4.5cm, inner sep=7pt] (m1) {No mide difusión, validación ni red completa.};
  \node[right=0.75cm of m1, draw=Warn, fill=SoftOrange, rounded corners, text width=4.5cm, inner sep=7pt] (m2) {Energía indirecta: sin medición eléctrica directa.};
  \node[below=0.45cm of m1, draw=Warn, fill=SoftOrange, rounded corners, text width=4.5cm, inner sep=7pt] (m3) {No modela incentivos ni participantes estratégicos.};
  \node[right=0.75cm of m3, draw=Warn, fill=SoftOrange, rounded corners, text width=4.5cm, inner sep=7pt] (m4) {Sin ASIC: no representa la frontera de Bitcoin.};
\end{tikzpicture}
\end{frame}

% 10
\begin{frame}{Conclusión: eficiencia no equivale a sostenibilidad}
\begin{columns}[T,totalwidth=\textwidth]
\column{0.55\textwidth}
\begin{itemize}\setlength\itemsep{0.28cm}
  \item El artículo es valioso como benchmark de ingeniería.
  \item No demuestra superioridad del protocolo democrático frente a Bitcoin.
  \item La sostenibilidad exige evaluar red completa, incentivos, emisiones y seguridad.
  \item En ASIC modernos, la alta densidad térmica abre una variable adicional: recuperación de calor residual.
\end{itemize}

\column{0.42\textwidth}
\centering
\begin{tikzpicture}[every node/.style={font=\small}]
  \node[draw=DeepBlue, fill=SoftBlue, rounded corners, minimum width=2.7cm, minimum height=1cm] (asic) at (0,1.2) {ASIC};
  \node[draw=Accent, fill=SoftOrange, rounded corners, minimum width=2.7cm, minimum height=1cm] (heat) at (0,-0.2) {Calor residual};
  \node[draw=Good, fill=Good!10, rounded corners, minimum width=2.7cm, minimum height=1cm] (use) at (0,-1.6) {Uso térmico};
  \draw[-{Latex[length=2.5mm]}, thick, Accent] (asic) -- node[right]{\scriptsize kW concentrados} (heat);
  \draw[-{Latex[length=2.5mm]}, thick, Good] (heat) -- node[right]{\scriptsize captura útil} (use);
  \draw[rounded corners, draw=DarkGray, dashed] (-1.7,-2.25) rectangle (1.7,1.85);
\end{tikzpicture}
\end{columns}
\vspace{0.1cm}
{\scriptsize Referencias: Nakamoto (2008); Paul, Sarkar y Mukherjee (2014); Lugo y Pedraza (2020).}
\end{frame}

\end{document}
